%=======================================================================
%  Rapport TP1 - Commande Avancee
%  ESIX Normandie - 5A Systemes Embarques (MeSE) - 2026/2027
%  Compilation : pdflatex rapport.tex (x2)
%=======================================================================
\documentclass[11pt,a4paper]{report}

\usepackage[utf8]{inputenc}
\usepackage[T1]{fontenc}
\usepackage[french]{babel}
\usepackage{lmodern}
\usepackage[a4paper,margin=2.5cm]{geometry}
\usepackage{amsmath,amssymb}
\usepackage{graphicx}
\usepackage{booktabs}
\usepackage{array}
\usepackage{caption}
\usepackage{subcaption}
\usepackage{listings}
\usepackage{xcolor}
\usepackage{fancyhdr}
\usepackage[hidelinks]{hyperref}

%--- Mise en page ------------------------------------------------------
\pagestyle{fancy}
\fancyhf{}
\fancyhead[L]{\small TP1 -- Commande adaptative d'un pendule}
\fancyhead[R]{\small ESIX 5A -- MeSE}
\fancyfoot[C]{\thepage}
\renewcommand{\headrulewidth}{0.4pt}

%--- Style des listings Matlab ----------------------------------------
\definecolor{codegray}{gray}{0.95}
\definecolor{codegreen}{rgb}{0,0.5,0}
\definecolor{codeblue}{rgb}{0,0,0.6}
\lstset{
  language=Matlab,
  basicstyle=\ttfamily\footnotesize,
  backgroundcolor=\color{codegray},
  commentstyle=\color{codegreen},
  keywordstyle=\color{codeblue},
  numbers=left, numberstyle=\tiny, numbersep=6pt,
  breaklines=true, frame=single, captionpos=b,
  literate={é}{{\'e}}1 {è}{{\`e}}1 {ê}{{\^e}}1 {à}{{\`a}}1 {ù}{{\`u}}1 {ô}{{\^o}}1 {î}{{\^i}}1 {ç}{{\c c}}1
}

%--- Macro figure : affiche l'image si elle existe, sinon un cadre -----
\newcommand{\figtp}[4][0.78]{%
  \begin{figure}[htbp]
    \centering
    \IfFileExists{#2}{\includegraphics[width=#1\textwidth]{#2}}%
      {\fbox{\parbox[c][5cm][c]{#1\textwidth}{\centering\itshape\small
        Figure a inserer : \texttt{\detokenize{#2}}}}}
    \caption{#3}
    \label{#4}
  \end{figure}}

\begin{document}

%=======================================================================
%  PAGE DE GARDE
%=======================================================================
\begin{titlepage}
\centering
{\large Université de Caen Normandie\par}
{\large École d'ingénieurs ESIX Normandie\par}
\vspace{0.3cm}
{\large Filière Mécatronique et Systèmes Embarqués (MeSE) -- 5\textsuperscript{e} année\par}
\vspace{2.2cm}

\rule{\textwidth}{1pt}\\[0.6cm]
{\Huge\bfseries Commande Avancée\par}
\vspace{0.5cm}
{\LARGE TP n\textsuperscript{o}1 : Commande adaptative\\[0.2cm] à modèle de référence d'un pendule\par}
\vspace{0.6cm}
\rule{\textwidth}{1pt}\\[2.5cm]

\begin{minipage}{0.45\textwidth}
\begin{flushleft}
\textbf{Rédigé par :}\\[0.2cm]
Abdelmoutalib \textsc{Douadi}
\end{flushleft}
\end{minipage}
\hfill
\begin{minipage}{0.45\textwidth}
\begin{flushright}
\textbf{Encadrant :}\\[0.2cm]
Pr. Fouad \textsc{Giri}
\end{flushright}
\end{minipage}

\vfill
{\large Année universitaire 2026 -- 2027\par}
\end{titlepage}

%=======================================================================
\tableofcontents
\clearpage

%=======================================================================
\chapter*{Introduction}
\addcontentsline{toc}{chapter}{Introduction}

Ce travail pratique porte sur la commande d'un pendule simple actionné par un couple moteur.
L'objectif pédagogique est de comparer, sur un même procédé non linéaire et sur les mêmes
scénarios d'essai, trois stratégies de commande de complexité croissante :

\begin{itemize}
  \item un régulateur \textbf{PID} à paramètres fixes, synthétisé par placement de pôles sur le
        modèle linéarisé autour de la position d'équilibre basse ;
  \item un régulateur \textbf{linéaire à modèle de référence} (MR), synthétisé sur le modèle
        discret ARMA du procédé, dont les paramètres sont supposés connus ;
  \item un régulateur \textbf{adaptatif à modèle de référence} (CAMR), qui estime en ligne les
        paramètres du régulateur par moindres carrés récursifs et n'exige donc aucune
        connaissance préalable du procédé.
\end{itemize}

Les trois régulateurs sont soumis aux mêmes perturbations : un changement d'amplitude de
consigne, qui éprouve la validité de l'hypothèse de linéarité, et un changement de la masse du
pendule, qui éprouve la robustesse vis-à-vis d'une variation paramétrique. Le rapport suit
l'ordre des tâches de l'énoncé et donne, pour chacune, le code produit, les résultats de
simulation et leur interprétation à la lumière du cours.

%=======================================================================
\chapter{Modélisation et simulation du procédé}

\section{Modèle de connaissance}

Le pendule est décrit par l'équation du mouvement
\begin{equation}
  J\,\ddot{\theta}(t) = C(t) - K\,\dot{\theta}(t) - m_0\,g\,r\,\sin\big(\theta(t)\big),
  \qquad J = m_0 r^2,
  \label{eq:pendule}
\end{equation}
où $\theta$ est l'angle par rapport à la verticale descendante, $C$ le couple de commande
(noté $u$ par la suite), $K$ le coefficient de frottement visqueux et $J$ le moment d'inertie.
Les valeurs numériques retenues sont rassemblées dans le tableau~\ref{tab:param}.

\begin{table}[htbp]
\centering
\caption{Paramètres du procédé et de la simulation}
\label{tab:param}
\begin{tabular}{llr}
\toprule
Grandeur & Symbole & Valeur \\
\midrule
Masse                       & $m_0$ & $0{,}25$ kg \\
Longueur du bras            & $r$   & $1$ m \\
Frottement visqueux         & $K$   & $0{,}1$ \\
Accélération de la pesanteur& $g$   & $10$ m/s$^2$ \\
Période d'échantillonnage   & $T_e$ & $0{,}05$ s \\
Pas de calcul (Runge--Kutta)& $T_c$ & $0{,}005$ -- $0{,}01$ s \\
\bottomrule
\end{tabular}
\end{table}

En posant le vecteur d'état $x = [\theta\;\; \dot{\theta}]^{T}$, l'équation~\eqref{eq:pendule}
s'écrit sous la forme d'état
\begin{equation}
  \dot{x} = f(x,u) =
  \begin{bmatrix} x_2 \\[2pt] \dfrac{u - K x_2 - m_0 g r \sin(x_1)}{m_0 r^2}\end{bmatrix}.
\end{equation}

\section{Simulation numérique : rôle de la méthode de Runge--Kutta}

Aucun banc réel n'étant disponible, le procédé est simulé numériquement. En raison du terme
$\sin(\theta)$, l'équation~\eqref{eq:pendule} n'admet pas de solution analytique : on l'intègre
donc pas à pas par la méthode de Runge--Kutta d'ordre 2 (méthode de Heun), qui prédit l'état
suivant par un pas d'Euler puis le corrige par la moyenne des deux pentes :
\begin{equation}
  x_1 = x(t) + T_c f\big(x(t),u\big), \qquad
  x(t+T_c) = x(t) + \frac{T_c}{2}\Big[f\big(x(t),u\big) + f\big(x_1,u\big)\Big].
\end{equation}

La structure à deux boucles imbriquées du script reproduit fidèlement une commande numérique :
la boucle interne effectue $m = T_e/T_c$ pas d'intégration et représente l'évolution
\emph{continue} du procédé ; la boucle externe, cadencée à $T_e$, représente le calculateur, qui
ne lit le capteur et ne recalcule la commande qu'une fois par période. Entre deux calculs, $u$
reste constante : c'est le \textbf{bloqueur d'ordre zéro}, dont on retrouvera l'effet lors de
l'échantillonnage de la fonction de transfert.

\section{Tâche 0 : comportement en boucle ouverte}

Le pendule est soumis à des échelons de couple d'amplitudes $0{,}1$, $0{,}5$, $1$ et
$1{,}5$~N$\cdot$m, sans aucune régulation.

\figtp{figures/BO_tache0_fig1.png}{Réponses en boucle ouverte à des échelons de couple}{fig:bo}

\paragraph{Résultats.} Les valeurs finales obtenues sont rassemblées dans le
tableau~\ref{tab:bo}, où elles sont comparées à la prédiction du modèle linéarisé
$\theta_\infty = u/(m_0 g r)$.

\begin{table}[htbp]
\centering
\caption{Position d'équilibre en boucle ouverte : procédé réel et modèle linéarisé}
\label{tab:bo}
\begin{tabular}{cccc}
\toprule
Échelon (N$\cdot$m) & $\theta_\infty$ réel (rad) & $\theta_\infty$ linéarisé (rad) & Écart relatif \\
\midrule
0,1 & 0,040 & 0,040 & $<1\%$ \\
0,5 & 0,201 & 0,200 & $0{,}5\%$ \\
1,0 & 0,412 & 0,400 & $3\%$ \\
1,5 & 0,644 & 0,600 & $7\%$ \\
\bottomrule
\end{tabular}
\end{table}

\paragraph{Interprétation.} Trois observations structurent la suite du travail.

\textbf{(i) Le procédé est très peu amorti.} Le modèle linéarisé donne un facteur
d'amortissement $\xi = a_1/(2\sqrt{a_0}) = 0{,}4/(2\sqrt{10}) \approx 0{,}063$ et une pulsation
propre $\omega_n = \sqrt{10} \approx 3{,}16$~rad/s. Les réponses présentent donc des
oscillations de période voisine de $2$~s, qui ne s'éteignent qu'au bout d'une vingtaine de
secondes ($\tau = 1/0{,}2 = 5$~s). Une régulation est indispensable pour amortir ce
comportement.

\textbf{(ii) La non-linéarité se manifeste dès les amplitudes moyennes.} L'équilibre réel vérifie
$m_0 g r \sin\theta_\infty = u$, alors que le modèle linéarisé suppose
$m_0 g r\,\theta_\infty = u$. Comme $\sin\theta < \theta$, le pendule réel doit monter
\emph{plus haut} que ne le prédit le modèle pour équilibrer le même couple, d'où un écart
croissant, qui atteint $7\%$ à $1{,}5$~N$\cdot$m. La période des oscillations s'allonge également
avec l'amplitude, phénomène qu'aucun modèle linéaire ne peut reproduire.

\textbf{(iii) Le procédé admet une limite physique.} Pour $u > m_0 g r = 2{,}5$~N$\cdot$m,
l'équation $\sin\theta_\infty = u/(m_0 g r)$ n'a plus de solution : il n'existe plus de position
d'équilibre et le pendule effectue des tours complets.

%=======================================================================
\chapter{Commande par régulateur PID}

\section{Synthèse par placement de pôles}

Autour de la position d'équilibre $\theta = 0$, l'approximation $\sin\theta \simeq \theta$
conduit à la fonction de transfert
\begin{equation}
  G(p) = \frac{\Theta(p)}{U(p)} = \frac{b}{p^2 + a_1 p + a_0},
  \qquad
  b = \frac{1}{m_0 r^2} = 4,\quad a_1 = \frac{K}{m_0 r^2} = 0{,}4,\quad a_0 = \frac{g}{r} = 10 .
  \label{eq:G}
\end{equation}

Le régulateur PID
$u = K_p\big[e + \frac{1}{T_i}\int e\,\mathrm{d}t + T_d\,\dot{e}\big]$, avec $e = y^{\mathrm{ref}} - y$,
place les trois pôles de la boucle fermée en $-\alpha$. L'identification du polynôme
caractéristique avec $(p+\alpha)^3$ fournit les relations du cours :
\begin{equation}
  K_p = \frac{3\alpha^2 - a_0}{b}, \qquad
  T_i = \frac{b K_p}{\alpha^3}, \qquad
  T_d = \frac{3\alpha - a_1}{b K_p}.
  \label{eq:pid}
\end{equation}

\paragraph{Tâche 1a.} Pour $\alpha = 2$, l'application numérique donne
\begin{equation}
  \boxed{K_p = 0{,}5}\qquad \boxed{T_i = 0{,}25~\mathrm{s}}\qquad \boxed{T_d = 2{,}8~\mathrm{s}}
\end{equation}

\begin{lstlisting}[caption={Tâche 1a : calcul des paramètres du PID}]
alpha=2;
Kp=(3*alpha^2 - a0)/b;      % gain proportionnel
Ti=b*Kp/alpha^3;            % constante de temps integrale
Td=(3*alpha - a1)/(b*Kp);   % constante de temps derivee
\end{lstlisting}

\section{Discrétisation de la loi de commande}

Le calculateur ne dispose de la mesure qu'aux instants $kT_e$. Chaque terme est donc
approché :
\begin{equation}
  \int_0^{t} e\,\mathrm{d}\tau \;\simeq\; I(k) = I(k-1) + T_e\,e(k)
  \qquad\text{(méthode des rectangles)},
\end{equation}
\begin{equation}
  \frac{\mathrm{d}e}{\mathrm{d}t} \;\simeq\; \frac{e(k)-e(k-1)}{T_e}
  \qquad\text{(différence arrière)} .
\end{equation}
L'intégrale est ainsi accumulée récursivement : à chaque période, on ajoute l'aire d'un
rectangle de largeur $T_e$ et de hauteur $e(k)$. Cette mémorisation est précisément ce qui
annule l'erreur statique : lorsque $e = 0$, $I$ cesse d'évoluer mais conserve la valeur qui
fournit le couple de maintien $m_0 g r\sin\theta$.

\begin{lstlisting}[caption={Tâche 1b : loi de commande PID discrétisée}]
e  = yref - y;                            % erreur de poursuite
I  = I + Te*e;                            % integrale de (yref-y)
u  = Kp*( e + (1/Ti)*I + Td*(e-e1)/Te );  % loi de commande PID
e1 = e;                                   % memorisation pour le terme derive
\end{lstlisting}

\section{Tâche 1 : consigne d'amplitude 0,1 puis 5}

Le signal de consigne est un carré de période $30$~s, d'amplitude $0{,}1$ sur $[0,75]$~s puis
$5$ sur $[75,150]$~s.

\figtp{figures/PID_tache1_fig1.png}{Tâche 1 : consigne et sortie du système}{fig:pid1a}
\figtp{figures/PID_tache1_fig2.png}{Tâche 1 : signal de commande}{fig:pid1b}

\paragraph{Première moitié ($A = 0{,}1$).} La sortie rejoint la consigne en environ $4$~s,
\emph{sans erreur statique}, conformément au rôle de l'action intégrale. La réponse présente
toutefois une forme caractéristique : un bond initial, un léger creux, puis la montée finale.
Ce comportement est dû au terme dérivé : au front de consigne, $e$ subit un saut et le terme
$K_p T_d (e-e_1)/T_e$ atteint $28 \times 0{,}1 \approx 2{,}8$~N$\cdot$m pendant une seule
période. Ce phénomène, connu sous le nom de \emph{derivative kick}, est visible sous forme de
pics sur la figure~\ref{fig:pid1b}. En pratique, on l'évite en dérivant la mesure plutôt que
l'erreur.

\paragraph{Seconde moitié ($A = 5$).} Le premier créneau d'amplitude $5$ n'apparaît qu'à
$t = 90$~s, car à $t = 75$~s le carré est dans sa phase basse. La boucle reste stable, mais les
performances se dégradent nettement : dépassement jusqu'à $\theta \approx 6{,}3$~rad
(soit $26\%$, le pendule effectuant presque un tour complet), excursion jusqu'à $-1{,}3$~rad au
retour à zéro et commande atteignant $140$~N$\cdot$m, valeur irréaliste pour un actionneur.

\paragraph{Interprétation.} Le réglage~\eqref{eq:pid} repose sur la linéarisation autour de
$\theta = 0$. Autour d'un point de fonctionnement $\bar\theta$ quelconque, le terme de raideur
devient $g\cos(\bar\theta)/r$. Pour $\bar\theta = 5$~rad, on a $g\cos(5)/r \approx 2{,}84$ au
lieu de $10$ : les pôles de la boucle fermée quittent le triple pôle en $-2$ et se placent en
$-5{,}38$ et $-0{,}31 \pm 1{,}18\,\mathrm{j}$. Cette paire complexe faiblement amortie explique
le dépassement observé.

\section{Tâche 2 : consigne d'amplitude 0,1 puis 2,0}

\figtp{figures/PID_tache2_fig1.png}{Tâche 2 : consigne et sortie pour une amplitude finale de 2,0}{fig:pid2a}
\figtp{figures/PID_tache2_fig2.png}{Tâche 2 : signal de commande}{fig:pid2b}

\paragraph{Résultats.} Le comportement est, de façon contre-intuitive, \emph{plus dégradé}
qu'à l'amplitude $5$ : la sortie part vers $5$~rad, revient à $-1$~rad et ne se stabilise sur la
consigne qu'au bout d'une quinzaine de secondes. La commande culmine à $57$~N$\cdot$m.

\paragraph{Interprétation.} Autour de $\bar\theta = 2$~rad (environ $115^\circ$, donc au-dessus
de l'horizontale), la raideur locale vaut $g\cos(2)/r \approx -4{,}16$ : elle est
\emph{négative}, la gravité éloignant le pendule de ce point au lieu de l'y ramener. Le
polynôme caractéristique de la boucle fermée devient
\begin{equation}
  p^3 + \big(a_1 + b K_p T_d\big)p^2 + \big(g\cos\bar\theta/r + bK_p\big)p + \frac{bK_p}{T_i}
  = p^3 + 6p^2 - 2{,}16\,p + 8 ,
\end{equation}
dont un coefficient est négatif : la condition nécessaire de stabilité de Routh n'est pas
satisfaite. Les pôles valent $-6{,}52$ et $+0{,}26 \pm 1{,}08\,\mathrm{j}$ : la boucle est
\textbf{localement instable}, et le pendule ne se stabilise qu'après avoir quitté cette zone.
À l'amplitude $5$, au contraire, $\cos(5) > 0$, ce qui explique le résultat meilleur de la
tâche~1.

\section{Tâche 3 : variation de la masse}

La masse passe de $0{,}25$ à $0{,}5$~kg à $t = 150$~s, le régulateur restant inchangé.

\figtp{figures/PID_tache3_fig1.png}{Tâche 3 : effet d'un doublement de la masse}{fig:pid3a}
\figtp{figures/PID_tache3_fig2.png}{Tâche 3 : signal de commande}{fig:pid3b}

\paragraph{Résultats et interprétation.} La boucle reste stable et sans erreur statique, mais la
réponse devient environ deux fois plus lente : le temps de réponse passe d'environ $3{,}75$~s à
$7{,}7$~s. Doubler $m_0$ divise $b$ et $a_1$ par deux ($b = 2$, $a_1 = 0{,}2$) ; le polynôme
caractéristique devient $p^3 + 3p^2 + 11p + 4$, dont les racines sont $-0{,}40$ et
$-1{,}30 \pm 2{,}88\,\mathrm{j}$. Le pôle réel lent en $-0{,}40$ impose une constante de temps
de $2{,}5$~s et domine désormais la dynamique.

\paragraph{Conclusion partielle.} Un régulateur à paramètres fixes ne garantit ses performances
qu'au voisinage du point de fonctionnement et du jeu de paramètres ayant servi à sa synthèse.
Les deux faiblesses mises en évidence --- sensibilité à l'amplitude et sensibilité aux
variations paramétriques --- motivent respectivement le passage à une commande à modèle de
référence, puis à sa version adaptative.

%=======================================================================
\chapter{Modèle discret ARMA du pendule}

\section{Tâche 4 : fonction de transfert continue}

L'application numérique de~\eqref{eq:G} donne
\begin{equation}
  G(p) = \frac{4}{p^2 + 0{,}4\,p + 10}.
\end{equation}

\section{Tâche 5 : échantillonnage, pôles et zéros}

L'échantillonnage avec bloqueur d'ordre zéro à $T_e = 0{,}05$~s conduit à
\begin{equation}
  G_e(z) = \frac{0{,}004957\,z + 0{,}004924}{z^2 - 1{,}95550\,z + 0{,}98020}.
\end{equation}

\begin{table}[htbp]
\centering
\caption{Pôles et zéros des modèles continu et échantillonné}
\label{tab:pz}
\begin{tabular}{lcc}
\toprule
 & $G(p)$ & $G_e(z)$ \\
\midrule
Pôles & $-0{,}2 \pm 3{,}156\,\mathrm{j}$ & $0{,}9777 \pm 0{,}1556\,\mathrm{j}$ \\
Module / partie réelle & partie réelle $<0$ & $|z| = 0{,}990 < 1$ \\
Zéros & aucun & $-0{,}9934$ \\
\bottomrule
\end{tabular}
\end{table}

\paragraph{Interprétation.} On vérifie numériquement que
$\exp(T_e \times \text{pôles de } G) = \text{pôles de } G_e$, ce qui confirme la relation
$z = e^{pT_e}$ : le demi-plan gauche du plan $p$ est transformé en l'intérieur du cercle unité
du plan $z$, et les deux modèles sont stables.

Le point remarquable est l'apparition d'un \textbf{zéro d'échantillonnage} en $z = -0{,}9934$,
alors que $G(p)$ n'en possède aucun. Ce zéro est introduit par le bloqueur d'ordre zéro
appliqué à un système de degré relatif $2$. Il est situé à l'intérieur du cercle unité, donc le
modèle discret est à minimum de phase, ce qui autorise la synthèse d'un régulateur à modèle de
référence. Mais sa proximité avec $-1$ aura des conséquences très concrètes sur le signal de
commande (section~\ref{sec:t13}).

\section{Tâches 6 et 7 : réponses indicielles comparées}

\figtp{figures/MR_tache6_fig11.png}{Réponse à un échelon de 0,1 N$\cdot$m : procédé réel, $G(p)$ et $G_e(z)$}{fig:t6a}
\figtp{figures/MR_tache6_fig14.png}{Réponse à un échelon de 1,5 N$\cdot$m : l'écart au modèle linéaire devient visible}{fig:t6b}

\paragraph{Interprétation.} Deux conclusions distinctes se dégagent.

D'une part, $G(p)$ et $G_e(z)$ coïncident \emph{exactement} aux instants d'échantillonnage,
quelle que soit l'amplitude : c'est la propriété fondamentale de l'échantillonnage avec
bloqueur d'ordre zéro, exact pour une entrée en escalier. La discrétisation n'introduit donc
aucune erreur de modélisation pour ce type d'entrée.

D'autre part, le procédé réel s'écarte progressivement des deux modèles linéaires lorsque
l'amplitude croît, conformément au tableau~\ref{tab:bo}. L'erreur de modélisation ne provient
pas de l'échantillonnage, mais bien de la \textbf{linéarisation}.

\section{Tâche 8 : mise sous forme ARMA}

Matlab fournit $G_e$ sous forme d'un quotient de polynômes en $z$. En divisant numérateur et
dénominateur par $z^2$ :
\begin{equation}
  G_e(z) = \frac{b_0 z^{-1} + b_1 z^{-2}}{1 + a_1 z^{-1} + a_2 z^{-2}}
         = z^{-1}\,\frac{B(z^{-1})}{A(z^{-1})},
\end{equation}
d'où, en revenant au domaine temporel, la forme ARMA demandée :
\begin{equation}
  A(q^{-1})\,y(t) = q^{-d} B(q^{-1})\,u(t), \qquad d = 1,
  \label{eq:arma}
\end{equation}
avec
\begin{equation}
  A(q^{-1}) = 1 - 1{,}95550\,q^{-1} + 0{,}98020\,q^{-2},
  \qquad
  B(q^{-1}) = 0{,}004957 + 0{,}004924\,q^{-1}.
\end{equation}
Le retard $d = 1$ traduit la causalité du calculateur : la commande calculée à l'instant $t$ ne
peut agir sur la sortie qu'à partir de l'instant $t+1$.

%=======================================================================
\chapter{Spécification des performances}

\section{Tâche 9 : choix et interprétation du modèle de référence}

Le modèle de référence
\begin{equation}
  G_M(p) = \frac{1}{\dfrac{p^2}{\omega_0^2} + \dfrac{2\xi}{\omega_0}p + 1},
  \qquad \xi = 1,\quad \omega_0 = 1~\mathrm{rad/s},
\end{equation}
ne décrit pas le procédé : il décrit le \emph{comportement souhaité} de la boucle fermée. La
consigne traverse $G_M$ pour produire la trajectoire de référence $y^*(t)$ que la sortie devra
suivre. Les valeurs choisies se justifient ainsi.

\begin{itemize}
  \item \textbf{Gain statique unitaire} : $G_M(0) = 1$, donc la sortie tend vers la consigne en
        régime permanent, sans erreur statique.
  \item \textbf{$\xi = 1$ (amortissement critique)} : les deux pôles sont confondus en
        $p = -\omega_0$, la réponse indicielle est donc apériodique, sans dépassement. Ce choix
        est pertinent face à un procédé dont l'amortissement naturel n'est que de $0{,}063$.
  \item \textbf{$\omega_0 = 1$~rad/s} : le temps de réponse à $5\%$ vaut $\approx 4{,}75/\omega_0
        \approx 4{,}75$~s. La boucle fermée est donc volontairement \emph{plus lente} que la
        dynamique propre du pendule ($\omega_n \approx 3{,}16$~rad/s), ce qui limite l'amplitude
        de la commande : un modèle de référence trop rapide exigerait des couples irréalistes et
        saturerait l'actionneur.
  \item \textbf{Réalisabilité} : le degré relatif de $G_M$ (égal à $2$) est supérieur ou égal au
        retard $d = 1$ du procédé, condition nécessaire pour que la poursuite soit réalisable.
\end{itemize}

\section{Tâche 10 : modèle de référence discret}

\begin{equation}
  G_M(p) = \frac{1}{p^2 + 2p + 1}
  \qquad\Longrightarrow\qquad
  G_{Me}(z) = \frac{0{,}001209\,z + 0{,}001169}{z^2 - 1{,}90246\,z + 0{,}90484}.
\end{equation}
Les pôles de $G_{Me}$ sont doubles en $z = 0{,}95123 = e^{-1\times 0{,}05}$, ce qui confirme de
nouveau la relation $z = e^{pT_e}$.

\figtp[0.7]{figures/MR_tache10_fig10.png}{Réponse indicielle du modèle de référence, continu et échantillonné}{fig:gm}

\section{Tâche 11 : polynôme de régulation $C(q^{-1})$}

Le polynôme $C(q^{-1}) = 1 + c\,q^{-1}$ fixe la dynamique de l'erreur de régulation. Sa racine
est $z = -c$ ; la condition $|c| < 1$ garantit qu'elle est à l'intérieur du cercle unité, donc
que cette dynamique est stable. Deux valeurs sont étudiées :
\begin{itemize}
  \item $c = 0$ : $C = 1$, l'erreur est annulée en temps minimal (réponse \emph{dead-beat}) ;
  \item $c = -0{,}5$ : un pôle en $0{,}5$ ralentit volontairement la correction, ce qui adoucit
        le signal de commande.
\end{itemize}

%=======================================================================
\chapter{Régulateur linéaire à modèle de référence}

\section{Synthèse : équation de Bezout}

La loi de commande à modèle de référence s'écrit
\begin{equation}
  R(q^{-1})B(q^{-1})\,u(t) + S(q^{-1})\,y(t) = C(q^{-1})\,y^*(t+d),
  \label{eq:loiMR}
\end{equation}
où $R$ et $S$ sont solutions de l'équation de Bezout
\begin{equation}
  A(q^{-1})R(q^{-1}) + q^{-d} S(q^{-1}) = C(q^{-1}).
  \label{eq:bezout}
\end{equation}

\paragraph{Degrés des polynômes.} Le cours donne $\deg R = d - 1$ et
$\deg S = \max(n_A - 1,\, n_C - d)$. Ici $d = 1$, $n_A = 2$ et $n_C = 1$, d'où
$\deg R = 0$ et $\deg S = 1$, soit $S(q^{-1}) = s_0 + s_1 q^{-1}$.

\paragraph{Justification de $R = 1$.} En évaluant~\eqref{eq:bezout} en $q^{-1} = 0$, le terme
$q^{-1}S$ s'annule et il reste $A(0)\,R = C(0)$. Comme $A(0) = C(0) = 1$, on obtient $R = 1$ :
le polynôme $R$ est unitaire, ce qui correspond à la condition de causalité du régulateur.

\paragraph{Détermination de $s_0$ et $s_1$.} En reportant $R = 1$ dans~\eqref{eq:bezout} :
\begin{equation}
  \big(1 + a_1 q^{-1} + a_2 q^{-2}\big) + q^{-1}\big(s_0 + s_1 q^{-1}\big) = 1 + c\,q^{-1},
\end{equation}
soit, après regroupement,
\begin{equation}
  1 + (a_1 + s_0)q^{-1} + (a_2 + s_1)q^{-2} = 1 + c\,q^{-1}.
\end{equation}
L'identification terme à terme fournit
\begin{equation}
  \boxed{s_0 = c - a_1} \qquad\text{et}\qquad \boxed{s_1 = -a_2}.
\end{equation}

\section{Tâche 12 : implantation de la loi de commande}

Avec $R = 1$ et $d = 1$, la relation~\eqref{eq:loiMR} s'écrit explicitement
\begin{equation}
  b_0 u(t) + b_1 u(t-1) + s_0 y(t) + s_1 y(t-1) = y^*(t+1) + c\,y^*(t),
\end{equation}
d'où la forme programmée :
\begin{equation}
  u(t) = \frac{1}{b_0}\Big[-b_1 u(t-1) - s_0 y(t) - s_1 y(t-1) + y^*(t+1) + c\,y^*(t)\Big].
\end{equation}

\begin{lstlisting}[caption={Tâche 12 : loi de commande à modèle de référence}]
R=1;
s0=c-a1;  s1=-a2;
u = ( -b1*u1 - s0*y - s1*y1 + Yref(k+1) + c*Yref(k) )/b0;
\end{lstlisting}

\section{Tâche 13 : consigne d'amplitude 0,1 puis 5}
\label{sec:t13}

\figtp{figures/MR_tache13_c0_fig1.png}{Tâche 13 : consigne et signal de référence}{fig:t13a}
\figtp{figures/MR_tache13_c0_fig2.png}{Tâche 13 : sortie, référence et consigne}{fig:t13b}
\figtp{figures/MR_tache13_c0_fig3.png}{Tâche 13 : signal de commande ($c = 0$)}{fig:t13c}

\paragraph{Qualité de poursuite.} Elle est excellente : l'erreur $|y - y^*|$ reste inférieure à
$10^{-3}$~rad sur la première moitié, alors même que le pendule part de $\theta_0 = 1$~rad avec
une vitesse initiale non nulle. Avec $c = 0$, la sortie rejoint la trajectoire de référence en
une seule période d'échantillonnage, ce qui correspond bien à la dynamique d'erreur en temps
minimal imposée par $C = 1$. Sur la seconde moitié, à l'amplitude $5$, l'erreur ne dépasse pas
$0{,}15$~rad et la commande reste inférieure à $5$~N$\cdot$m, ce qui constitue déjà une
amélioration spectaculaire par rapport au PID de la tâche~1 ($6{,}3$~rad de dépassement et
$140$~N$\cdot$m).

\paragraph{Le prix caché : l'oscillation du signal de commande.} La figure~\ref{fig:t13c}
révèle un phénomène absent de la sortie : la commande change de signe à chaque période, avec un
pic initial d'environ $350$~N$\cdot$m, et met près de $60$~s à s'éteindre.

L'explication tient au zéro identifié à la tâche~5. La loi~\eqref{eq:loiMR} fait intervenir
$B(q^{-1})u(t)$ : le régulateur \emph{simplifie} le polynôme $B$, dont le zéro
$z = -b_1/b_0 = -0{,}9934$ devient un \textbf{pôle de la dynamique de la commande} :
\begin{equation}
  u(t) \simeq -\frac{b_1}{b_0}\,u(t-1) = (-0{,}9934)\,u(t-1).
\end{equation}
La commande oscille donc à la fréquence de Nyquist et décroît avec une constante de temps
$T_e/\ln(1/0{,}9934) \approx 7{,}5$~s. Ce mode n'apparaît pas sur $y$ puisque le zéro qui l'a
engendré le compense exactement en sortie : il s'agit d'une \emph{oscillation cachée}, qui
n'en solliciterait pas moins l'actionneur réel. C'est la limite pratique de l'hypothèse de
minimum de phase lorsqu'un zéro est trop proche du cercle unité --- situation d'autant plus
fréquente que $T_e$ est petit.

\paragraph{Effet du polynôme $C$.} Avec $c = -0{,}5$, le pic initial tombe à environ
$125$~N$\cdot$m et l'extinction est plus rapide. Le prix à payer est une convergence un peu plus
lente de l'erreur de poursuite : $C$ réalise donc un compromis explicite entre rapidité et
douceur de la commande.

\section{Tâche 14 : consigne d'amplitude 0,1 puis 2,0}

\figtp{figures/MR_tache14_c0_fig2.png}{Tâche 14 : sortie et référence pour une amplitude finale de 2,0}{fig:t14}

\paragraph{Résultats et interprétation.} L'erreur de poursuite reste de l'ordre de
$0{,}03$~rad et la commande n'excède pas $2{,}5$~N$\cdot$m. Là où le PID devenait
\emph{localement instable} (tâche~2), le régulateur à modèle de référence reste stable et
précis. La raison est structurelle : il ne repose pas sur un placement de pôles figé autour
d'un point d'équilibre, mais sur l'inversion du modèle à chaque période, le terme
$S(q^{-1})y(t)$ compensant en permanence le couple de gravité, assimilable à une perturbation
lentement variable. L'erreur résiduelle, non nulle, provient de ce que le modèle ARMA utilisé
reste linéarisé autour de $\theta = 0$.

\section{Tâche 15 : variation de la masse}

\figtp{figures/MR_tache15_c0_fig2.png}{Tâche 15 : sortie et référence après doublement de la masse}{fig:t15a}
\figtp{figures/MR_tache15_c0_fig3.png}{Tâche 15 : réapparition de l'oscillation de commande après $t = 75$~s}{fig:t15b}

\paragraph{Résultats et interprétation.} La poursuite reste bonne, mais l'erreur est multipliée
par un facteur voisin de $8$ après le changement, et l'oscillation alternée de la commande
réapparaît puis s'éteint lentement. Le régulateur continue en effet d'utiliser les coefficients
$b_0$, $b_1$, $s_0$ et $s_1$ calculés pour le système initial : la simplification du polynôme
$B$ n'est plus exacte, et la dynamique d'erreur n'est plus celle imposée par $C$.

La dégradation reste ici modérée, car l'écart paramétrique est modeste et connu. Dans une
application réelle --- un bras portant une charge inconnue et variable --- l'écart peut être
arbitrairement grand. D'où la nécessité d'estimer les paramètres \emph{en ligne}.

%=======================================================================
\chapter{Commande adaptative à modèle de référence}

\section{Principe et paramétrisation}

Les paramètres du modèle~\eqref{eq:arma} étant désormais supposés inconnus, les coefficients du
régulateur ne peuvent plus être calculés hors ligne. On réalise donc une version adaptative du
régulateur précédent, suivant la démarche du chapitre~3 du cours. Avec
$\boldsymbol{R}(q^{-1}) = R(q^{-1})B(q^{-1}) = B(q^{-1})$, le vecteur des paramètres du
régulateur et le vecteur de régression se réduisent à
\begin{equation}
  \Theta = \big[\boldsymbol{r}_0,\ \boldsymbol{r}_1,\ s_0,\ s_1\big]^{T} \in \mathbb{R}^4,
  \qquad
  \phi(t) = \big[u(t-1),\ u(t-2),\ y(t-1),\ y(t-2)\big]^{T} \in \mathbb{R}^4 ,
\end{equation}
puisque $\boldsymbol{r}_2 = 0$. Cette simplification réduit l'exigence d'excitation
persistante.

\section{Tâche 16 : estimateur paramétrique et validation en boucle ouverte}

L'estimateur des moindres carrés récursifs s'écrit
\begin{align}
  \epsilon(t)      &= C(q^{-1})y(t) - \phi^{T}(t)\,\hat\Theta(t-1), \\
  \hat\Theta(t)    &= \hat\Theta(t-1) + \frac{P(t-1)\phi(t)\,\epsilon(t)}{1 + \phi^{T}(t)P(t-1)\phi(t)},
  \label{eq:rls1}\\
  P(t)             &= P(t-1) - \frac{P(t-1)\phi(t)\phi^{T}(t)P(t-1)}{1 + \phi^{T}(t)P(t-1)\phi(t)} .
  \label{eq:rls2}
\end{align}

\paragraph{Remarque sur l'ordre des mises à jour.} L'énoncé fait apparaître $P(t)$ au numérateur
de~\eqref{eq:rls1} et $P(t-1)$ au dénominateur. Les deux écritures sont équivalentes, car
en multipliant~\eqref{eq:rls2} par $\phi(t)$ on obtient l'identité
\begin{equation}
  P(t)\phi(t) = \frac{P(t-1)\phi(t)}{1 + \phi^{T}(t)P(t-1)\phi(t)} .
\end{equation}
Le facteur de normalisation est donc \emph{déjà contenu} dans $P(t)$. En pratique, on met à jour
$\hat\Theta$ avant $P$ en conservant le dénominateur (version programmée ci-dessous), ou bien
$P$ avant $\hat\Theta$ en le supprimant. Appliquer le dénominateur \emph{et} utiliser $P(t)$
reviendrait à normaliser deux fois et ralentirait notablement la convergence.

\begin{lstlisting}[caption={Tâche 16 : estimateur des moindres carrés récursifs}]
epsilon = (y + c*y1) - phi'*Theta;
Theta   = Theta + P*phi*epsilon/( 1 + phi'*P*phi );
P       = P - (P*phi*phi'*P)/( 1 + phi'*P*phi );
\end{lstlisting}

\paragraph{Validation en boucle ouverte.} L'estimateur est d'abord testé seul, le procédé étant
excité par le signal carré de consigne ($u(t) = $ consigne).

\figtp{figures/CAMR_tache16_fig6.png}{Tâche 16 : estimation de $b_0$ et $b_1$ en boucle ouverte}{fig:t16a}
\figtp{figures/CAMR_tache16_fig7.png}{Tâche 16 : estimation de $s_0$ et $s_1$ en boucle ouverte}{fig:t16b}

\begin{table}[htbp]
\centering
\caption{Tâche 16 : paramètres estimés en boucle ouverte et valeurs vraies}
\label{tab:t16}
\begin{tabular}{lcc}
\toprule
Paramètre & Estimé & Valeur vraie \\
\midrule
$b_0$ & $\approx 0{,}0050$   & $0{,}004957$ \\
$b_1$ & $\approx 0{,}0047$   & $0{,}004924$ \\
$s_0$ & $\approx 1{,}9566$   & $1{,}95550$ \\
$s_1$ & $\approx -0{,}9805$  & $-0{,}98020$ \\
\bottomrule
\end{tabular}
\end{table}

\paragraph{Interprétation.} Les estimés convergent en $2$ à $3$~secondes et l'erreur
d'estimation $\epsilon$ chute très rapidement vers zéro. Le signal carré constitue donc une
\textbf{excitation persistante} suffisante : ses fronts contiennent assez de composantes
fréquentielles pour identifier les quatre paramètres. L'écart résiduel, de l'ordre du pourcent,
s'explique par la non-linéarité du procédé, que le modèle ARMA ne peut reproduire exactement :
l'estimateur converge en réalité vers le \emph{meilleur modèle linéaire} au sens des moindres
carrés pour l'excitation appliquée.

\section{Tâche 17 : loi de commande adaptative}

La loi de commande~\eqref{eq:loiMR} est reprise en remplaçant les paramètres vrais par leurs
estimés courants (principe d'équivalence certaine) :
\begin{equation}
  u(t) = \frac{1}{\hat{\boldsymbol{r}}_0(t)}
  \Big[-\hat{\boldsymbol{r}}_1(t)u(t-1) - \hat{s}_0(t)y(t) - \hat{s}_1(t)y(t-1)
  + y^*(t+1) + c\,y^*(t)\Big].
\end{equation}

\begin{lstlisting}[caption={Tâche 17 : loi de commande adaptative}]
beta0=Theta(1);
if (abs(Theta(1))<0.001)
    beta0=(0.001)*sign(Theta(1));   % protection : r0 estime ne doit pas s'annuler
end
u = -Theta(2)*u1 - Theta(3)*y - Theta(4)*y1 + Yref(k+1) + c*Yref(k);
u = u/beta0;
\end{lstlisting}

La protection sur $\hat{\boldsymbol{r}}_0$ est indispensable : la division par un estimé proche
de zéro ferait diverger la commande pendant la phase transitoire d'apprentissage.

\section{Tâche 18 : consigne d'amplitude 0,1 puis 5}

\figtp{figures/CAMR_tache18_c0_fig2.png}{Tâche 18 : sortie, référence et consigne}{fig:t18a}
\figtp{figures/CAMR_tache18_c0_fig3.png}{Tâche 18 : signal de commande}{fig:t18b}
\figtp{figures/CAMR_tache18_c0_fig5.png}{Tâche 18 : erreur d'estimation $\epsilon$}{fig:t18c}

\paragraph{Résultats.} Les estimés convergent en moins d'une seconde, soit bien plus vite qu'en
boucle ouverte : la condition initiale $\theta_0 = 1$~rad et l'action du régulateur excitent
fortement le procédé dès les premiers instants. L'erreur de poursuite s'établit ensuite autour
de $4\cdot10^{-4}$~rad, puis d'environ $0{,}09$~rad après le passage à l'amplitude $5$.

\paragraph{Comparaison avec la tâche 13.} La qualité de poursuite est au moins équivalente à
celle du régulateur linéaire, alors que le régulateur adaptatif part de
$\hat\Theta(0) = [0{,}01\ 0\ 0\ 0]^T$, c'est-à-dire d'une méconnaissance totale du procédé.
C'est le résultat central de cette partie. Le signal de commande est en revanche plus agité aux
grandes amplitudes : les figures des paramètres montrent que $\hat b_0$ dérive vers $0{,}0037$
au lieu de $0{,}00496$. Ce n'est pas un défaut de l'estimateur mais une conséquence logique du
fait qu'à $5$~rad le procédé n'est plus décrit par le modèle linéarisé autour de $0$ :
l'estimateur suit un modèle linéaire local qui change avec le point de fonctionnement.

\section{Tâche 19 : consigne d'amplitude 0,1 puis 2,0}

\figtp{figures/CAMR_tache19_c0_fig2.png}{Tâche 19 : sortie et référence pour une amplitude finale de 2,0}{fig:t19}

\paragraph{Comparaison avec la tâche 14.} Les deux régulateurs donnent la même erreur de
poursuite, de l'ordre de $0{,}027$~rad. L'adaptation \emph{n'apporte rien} dans ce scénario, et
ce résultat est instructif : le procédé n'a pas changé, seule la consigne a changé. L'erreur
résiduelle ne provient pas d'une méconnaissance des paramètres, mais d'une \emph{erreur de
structure} du modèle (linéarisation d'un procédé non linéaire), qu'aucun des deux régulateurs
ne peut annuler. On retiendra que l'adaptation compense une variation du procédé, non une
inadéquation de la structure du modèle.

Il convient néanmoins de rappeler qu'à cette amplitude le PID de la tâche~2 devenait localement
instable, alors que les deux régulateurs à modèle de référence restent stables.

\section{Tâche 20 : variation de la masse}

\figtp{figures/CAMR_tache20_c0_fig6.png}{Tâche 20 : poursuite des nouvelles valeurs de $b_0$ et $b_1$ après le changement de masse}{fig:t20a}
\figtp{figures/CAMR_tache20_c0_fig2.png}{Tâche 20 : sortie et référence}{fig:t20b}

\paragraph{Résultats.} À $t = 100$~s, les valeurs vraies subissent un saut ($b_0$ passe de
$0{,}00496$ à $0{,}00249$). Les estimés les rejoignent en moins d'une seconde, et l'erreur de
poursuite revient à son niveau nominal, de l'ordre de $2\cdot10^{-3}$~rad.

\paragraph{Comparaison avec la tâche 15.} C'est dans ce scénario que la commande adaptative
démontre sa supériorité. Le régulateur linéaire conservait des paramètres erronés, voyait son
erreur multipliée par $8$ et son signal de commande se dégrader durablement. Le régulateur
adaptatif, lui, \emph{recalcule ses propres coefficients} et retrouve les performances fixées
par le modèle de référence, sans qu'aucune information sur la nouvelle masse ne lui ait été
fournie.

\paragraph{Rôle de la réinitialisation de $P$.} Si l'on supprime l'instruction
\texttt{P=1000*eye(4)} du bloc de changement, les estimés ne bougent pratiquement plus après
$t = 100$~s ($\hat b_0$ reste à $0{,}0049$ au lieu de descendre à $0{,}0025$). La matrice $P$
étant décroissante par construction~\eqref{eq:rls2}, elle devient quasi nulle après une longue
phase de convergence : le gain d'adaptation s'annule et l'estimateur « s'endort ». C'est le
phénomène d'\textbf{extinction du gain}. Deux remèdes classiques existent : la réinitialisation
de $P$, employée ici, et le \textbf{facteur d'oubli} $\lambda \in\,]0,1[$, qui consiste à
diviser $P$ par $\lambda$ à chaque période afin de maintenir l'estimateur en éveil de façon
permanente. Ce dernier est préférable en pratique, puisqu'il ne suppose pas de connaître
l'instant de la variation.

%=======================================================================
\chapter*{Synthèse et conclusion}
\addcontentsline{toc}{chapter}{Synthèse et conclusion}

Le tableau~\ref{tab:synthese} résume le comportement des trois régulateurs face aux trois
scénarios d'essai.

\begin{table}[htbp]
\centering
\caption{Synthèse comparative des trois stratégies de commande}
\label{tab:synthese}
\begin{tabular}{p{3.2cm}p{3.4cm}p{3.4cm}p{3.4cm}}
\toprule
 & \textbf{PID} & \textbf{MR linéaire} & \textbf{MR adaptatif} \\
\midrule
Connaissance requise & modèle linéarisé & modèle ARMA complet & structure seule ($n$, $d$) \\
Petite amplitude & bon, $t_r \approx 4$~s & excellent, $|e|<10^{-3}$ & excellent après apprentissage \\
Amplitude 2,0 rad & \textbf{localement instable} & stable, $|e|\approx 0{,}03$ & stable, $|e|\approx 0{,}027$ \\
Amplitude 5 rad & dépassement $26\%$, $140$~N$\cdot$m & $|e|\approx 0{,}15$ & $|e|\approx 0{,}09$ \\
Masse doublée & deux fois plus lent & erreur $\times 8$ & performances conservées \\
Signal de commande & pics dérivés & oscillation cachée & agité aux fortes amplitudes \\
\bottomrule
\end{tabular}
\end{table}

Ce TP illustre une progression logique. Le PID, synthétisé par placement de pôles sur un modèle
linéarisé, n'offre les performances attendues qu'au voisinage immédiat du point de
fonctionnement : il perd sa stabilité lorsque la raideur locale change de signe, et sa rapidité
lorsque la masse varie. Le régulateur linéaire à modèle de référence, fondé sur l'inversion du
modèle discret, améliore considérablement la poursuite et la robustesse à l'amplitude, mais
suppose les paramètres connus et exacts ; de plus, la simplification du polynôme $B$ transfère
sur la commande un mode oscillant hérité du zéro d'échantillonnage, point à surveiller
impérativement dans toute implantation réelle. Enfin, la commande adaptative conserve les
performances du modèle de référence sans aucune connaissance préalable des paramètres et
s'ajuste automatiquement à leurs variations, à condition que l'excitation soit persistante et
que le gain d'adaptation soit maintenu, par réinitialisation de $P$ ou par facteur d'oubli.

Une limite commune aux trois approches mérite d'être soulignée : aucune ne supprime l'erreur
due à la \emph{structure} linéaire du modèle. La poursuite d'une consigne de grande amplitude
sur un procédé fortement non linéaire relèverait de techniques complémentaires, telles que la
linéarisation par bouclage ou le séquencement de gains, qui prolongent naturellement ce travail.

\end{document}
